\subsection{Ensemble MC: Almost Every Starting Point}
\label{sec:ensemble-variant}

The variants above change the \emph{selection rule}; here we change the
\emph{starting point}.  The generalized EM construction
(Section~\ref{sec:ensemble}) parametrizes the sequence by its initial
value~$n$: $\Prod{n}(0) = n$,
$\seq{n}(k) = \minFac(\Prod{n}(k)+1)$,
$\Prod{n}(k{+}1) = \Prod{n}(k)\cdot\seq{n}(k)$.
The standard orbit is $n = 2$.

\begin{definition}[\lean{EM/Ensemble/WeylChain.lean}{GenMullinConjecture}]
$\mathrm{GenMC}(n)$: every prime $q \nmid n$ appears in the generalized EM
sequence from~$n$, i.e.\ $\forall q$ prime with $q \nmid n$, $\exists k$,
$\seq{n}(k) = q$.  (Primes dividing~$n$ divide every $\Prod{n}(k)$ and are
never selected, so the unrestricted version is false for all $n \ge 2$;
dead end~\#175.)
\end{definition}

The reduction
\lean{EM/Ensemble/WeylChain.lean}{gen\_hitting\_implies\_gen\_mc\_proved}
{\color{leanink}$\checkmark$} proves that cofinal walk hitting
implies~$\mathrm{GenMC}(n)$ for every squarefree~$n$, mirroring the
standard inductive bootstrap.  And
\lean{EM/Ensemble/WeylChain.lean}{gen\_mc\_two\_implies\_mc}
{\color{leanink}$\checkmark$} shows $\mathrm{GenMC}(2) \Rightarrow
\mathrm{MC}$ (and \lean{EM/Ensemble/WeylChain.lean}{gen\_mc\_two\_iff\_mc}
{\color{leanink}$\checkmark$} the converse), so the generalized
conjecture specializes to the standard one.

\paragraph{Bag arithmetic.}
At each step~$k$, the Euclid number $\Prod{n}(k)+1$ has a ``bag'' of
prime factors; the selected prime $\seq{n}(k)$ is the minimum of this
bag.  The formalization
(\code{EM/Ensemble/BagArithmetic.lean}, 224~lines, zero \code{sorry})
develops the bag-level quantities:
\defname{genEuclidOmega}~$\omega(n,k) = |\text{primeFactors}(\Prod{n}(k)+1)|$
(at least~1,
\lean{EM/Ensemble/BagArithmetic.lean}{genEuclidOmega\_pos}
{\color{leanink}$\checkmark$}),
\defname{genBagDiversity} (number of distinct residue classes mod~$q$
among the factors, $\leq \omega$ and $\leq q$),
\defname{genFactorsInClass}~$F_a$ (factors in class~$a$, with the
partition identity
$\sum_a |F_a| = \omega$,
\lean{EM/Ensemble/BagArithmetic.lean}{genFactorsInClass\_card\_sum}
{\color{leanink}$\checkmark$}),
and \defname{genEuclidCofactor} (the quotient
$(\Prod{n}(k)+1)/\seq{n}(k)$, whose prime factors are a subset of
the bag,
\lean{EM/Ensemble/BagArithmetic.lean}{cofactor\_primeFactors\_subset}
{\color{leanink}$\checkmark$}).
For $k \geq 1$, all bag elements are $\geq 3$
(\lean{EM/Ensemble/BagArithmetic.lean}{primeFactors\_ge\_three\_of\_pos}
{\color{leanink}$\checkmark$}).

\paragraph{Density variants.}
Instead of proving $\mathrm{GenMC}(n)$ for the single orbit $n = 2$,
one can ask how many starting points satisfy it.  Two density levels
are formalized:

\begin{itemize}
\item \textbf{PositiveDensityRSD}
  (\lean{EM/Population/ReciprocalSum.lean}{PositiveDensityRSD}):
  a positive density~$\delta > 0$ of squarefree~$n$ have divergent
  reciprocal sums $\sum_k 1/\seq{n}(k) = \infty$.

\item \textbf{AlmostAllSquarefreeRSD}
  (\lean{EM/Population/ReciprocalSum.lean}{AlmostAllSquarefreeRSD}):
  density~$1$ of squarefree~$n$ have divergent reciprocal sums.
\end{itemize}

\noindent
The reciprocal sum is a natural proxy: if $\sum 1/\seq{n}(k)$
diverges, then the primes appearing in the sequence from~$n$ cannot be
confined to a thin set; although divergence alone does not formally
imply~$\mathrm{GenMC}(n)$, it rules out the worst obstruction (a
bounded sequence of captured primes).

\paragraph{The one-hypothesis route: LMG.}
The proved chain
\archived{EM/Archive/Population/ReciprocalSumArchive.lean}{lmg\_implies\_positive\_density\_rsd}
{\color{leanink}$\checkmark$} shows:
\[
  \mathrm{LMG} \;\Longrightarrow\; \mathrm{PositiveDensityRSD},
\]
where \textbf{LinearMeanGrowth}~(LMG) asks that
$\mathbb{E}_n[S_K(n)] \geq \kappa K$ for some $\kappa > 0$ and all
$K$ large enough (here $S_K(n) = \sum_{k<K} 1/\seq{n}(k)$ is the
partial reciprocal sum, averaged over squarefree~$n$).
LMG is a \emph{population-level} statement: it involves no specific
orbit, bypassing the orbit-specificity barrier that obstructs CME\@.
It is also strictly weaker than CME---it requires only that the
ensemble mean grows linearly, not conditional equidistribution.

With the additional hypothesis PairwiseStepDecorrelation~(PSD)---that
$\operatorname{Cov}[1/\seq{n}(j),\, 1/\seq{n}(k)] \to 0$ for
$j \neq k$---the stronger conclusion follows:
$\mathrm{PSD} + \mathrm{LMG} \Rightarrow \mathrm{AlmostAllSquarefreeRSD}$
(\archived{EM/Archive/Population/ReciprocalSumArchive.lean}{dead\_end\_125\_pairwise\_revival}
{\color{leanink}$\checkmark$}).
This is the revival of Dead End~\#125: pairwise cancellation does not
imply $k$-wise cancellation, but variance bounds only use $k = 2$, so
pairwise decorrelation is exactly enough.
The bridges are all machine-verified:
$\mathrm{IVB}(1/4)$ {\color{leanink}$\checkmark$},
$\mathrm{PSD} + \mathrm{IVB} \Rightarrow$ variance bound
{\color{leanink}$\checkmark$},
Chebyshev concentration {\color{leanink}$\checkmark$}.

\paragraph{The first moment, and a dead CRT route.}
LMG follows from the open hypothesis
\textbf{FirstMomentStep}~(FMS):
$\mathbb{E}_n[1/\seq{n}(k)] \to \kappa$ as $X \to \infty$ for each~$k$
(\lean{EM/Ensemble/FirstMoment.lean}{FirstMomentStep}).
The CRT-propagation and backward-dynamics route to FMS---via
CRTPropagationStep, EnsembleTransitionApprox, and Tao-style backward
density decomposition---is dead: the absorption mechanism falsifies
AccumulatorEquidistPropagation (Dead End~\#137), and the route is
catalogued, with its machine-verified but now vacuous reductions, in
Appendix~\ref{app:dead-ends}.

\paragraph{The accumulator as a commutative state.}
The EM accumulator $\Prod{n}(k) = \prod S_k$ is a squarefree integer
determined by its prime set~$S_k$.  The dynamics
$S \mapsto S \cup \{\minFac(\prod S + 1)\}$ is a well-defined map on
the lattice of finite prime sets, and the walk mod~$q$ depends on~$S$
only through $\prod S \bmod q$.  Two orbits from different starting
points that ever reach the same product mod~$q$ have approximately
identical future walks---the construction history is forgotten.

This approximate Markov property, combined with the super-exponential
growth $\Prod{n}(k) \geq 2^{2^k}$ (which makes CRT blindness
increasingly accurate), suggests that the walk mod~$q$ is effectively
ergodic from generic starting points.  Making this rigorous---showing
that the CRT error is summable along the orbit---would give a.a.\
$\mathrm{GenMC}$ for the standard sequence, without any modification
of the selection rule.

\paragraph{Comparison with the orbit hypothesis.}
CME (Section~\ref{sec:ensemble}) is a statement about the single
orbit $n = 2$.
The ensemble route here proves a density result for \emph{almost all}
starting points, contingent on CRTPropagationStep, which is a
population-level hypothesis with no orbit-specificity.
Neither subsumes the other: CME gives MC directly; CRTPropagationStep
gives PositiveDensityRSD (which does not formally imply MC for $n = 2$).
However, the ensemble route faces a structurally easier barrier, since
its open hypothesis is a statement about averages, not about one
specific orbit.

\subsection{Weak Mullin and Analytic Weakenings}
\label{sec:weak-mullin}

The variants above change the \emph{selection rule}; this subsection
changes the \emph{conclusion}.  Instead of proving that every prime
appears ($M = \emptyset$), one can ask that the set of missing
primes is small, or that the EM reciprocal sum diverges.

\paragraph{Missing primes and Weak Mullin.}
Define the set of \emph{missing primes}
\[
M = \{p \in \mathbb{P} : p \neq \seq{2}(k) \text{ for all } k\}.
\]
MC is the assertion $M = \emptyset$.  A natural weakening asks only that
$M$ is \emph{small} in an analytic sense:

\begin{property}[\defname{WeakMullin} (\abbr{WM})]
\lean{EM/Population/WeakMullin.lean}{WeakMullin}
The reciprocal sum over missing primes converges:
$\sum_{q \in M} 1/q < \infty$.
\end{property}

This is strictly weaker than MC (which gives $M = \emptyset$ and hence
a trivially convergent empty sum) but strong enough to have a nontrivial
consequence:

\begin{property}[\defname{ReciprocalDivergence} (\abbr{RD})]
\lean{EM/Population/WeakMullin.lean}{ReciprocalDivergence}
The reciprocal sum over EM primes diverges:
$\sum_{k=0}^{\infty} 1/\seq{2}(k) = \infty$.
\end{property}

\begin{theorem}[\lean{EM/Population/WeakMullin.lean}{mc\_implies\_wm}]
$\MC \Longrightarrow \mathrm{WM}$.
\end{theorem}

\begin{theorem}[\lean{EM/Population/WeakMullin.lean}{wm\_implies\_rd}]
\label{thm:wm-rd}
$\mathrm{WM} \Longrightarrow \mathrm{RD}$.
\end{theorem}

\noindent
The proof partitions $\mathbb{P} = A \cup M$ (appeared vs.\ missing),
notes that both $\Sigma_A$ and $\Sigma_M$ converge under WM $+$ $\neg$RD,
contradicting the divergence of $\sum_{p} 1/p$.

\paragraph{Product divergence.}
The negation of WM has a quantitative consequence: $\neg\mathrm{WM}$ gives
$\sum_{q \in M} 1/(q{-}1) = \infty$, so the ``safe fraction'' product
$\prod_{q \in F}(q{-}2)/(q{-}1) \leq e^{-\sum 1/(q-1)} \to 0$ and its
reciprocal $\prod (q{-}1)/(q{-}2)$ exceeds any bound
(\lean{EM/Population/WeakMullin.lean}{not\_wm\_product\_diverges}).
The product $\prod (q{-}1)/(q{-}2)$ over missing primes measures the
``confinement cost'': each missing prime~$q$ confines the walk to $q-2$
of the $q-1$ units in $(\ZZ/q\ZZ)^\times$ (avoiding~$-1$).  If WM
fails, this confinement cost is unbounded.

\paragraph{Confined energy.}
For a missing prime $q \geq 3$, the walk avoids $-1$ entirely, so the
$N$ visits are distributed among only $q-2$ residues.  A pigeonhole
argument gives a confined-energy lower bound
(\lean{EM/Population/WeakMullin.lean}{per\_prime\_confined\_energy\_lb}):
$N^2/(q{-}2) \leq (q{-}1)\sum_a V_N(a)^2$.
The excess energy ratio $(q{-}1)/(q{-}2)$ compared to the
equidistributed baseline $N^2/(q{-}1)$ is bounded away from~$1$.
\defname{JointSVE}
(\lean{EM/Population/WeakMullin.lean}{JointSVE})---joint subquadratic
visit energy over any finite set of missing primes---implies MMCSB
(\lean{EM/Population/WeakMullin.lean}{jointSVE\_implies\_mmcsb}),
connecting the ``missing primes'' picture to the spectral route.

\paragraph{Avoidance tube and tube collapse.}
Define the \emph{tube ratio} for a finite set~$F$ of primes $\geq 3$:
\[
  \tau(F) \;=\; \prod_{q \in F} \frac{q-2}{q-1}.
\]
This measures the fraction of $(\ZZ/q\ZZ)^\times$ classes that a walk
avoiding $-1$ at each~$q$ is allowed to visit.  The tube ratio is
monotone decreasing in~$F$ and satisfies $0 \leq \tau(F) \leq 1$
(\lean{EM/Population/AvoidanceTube.lean}{tubeRatio\_monotone}
{\color{leanink}$\checkmark$},
\lean{EM/Population/AvoidanceTube.lean}{tubeRatio\_le\_one}
{\color{leanink}$\checkmark$}).

\begin{theorem}[Tube collapse ---
\lean{EM/Population/AvoidanceTube.lean}{tube\_collapse}
{\color{leanink}$\checkmark$}]
\label{thm:tube-collapse}
$\neg\mathrm{WM} \;\Longrightarrow\;$ for every $\varepsilon > 0$, there
exists a finite set~$F \subseteq M$ with $\tau(F) < \varepsilon$.
\end{theorem}

\noindent
The proof uses the inequality
$\tau(F) \leq \exp(-\sum_{q \in F} 1/(q{-}1))$
together with the divergence of $\sum_{q \in M} 1/(q{-}1)$ under~$\neg$WM\@.
For any unit~$a$ with $V_N(a) = 0$, the confined energy gives a rogue
character:

\begin{theorem}[Rogue character ---
\lean{EM/Population/AvoidanceTube.lean}{rogue\_character\_exists}
{\color{leanink}$\checkmark$}]
\label{thm:rogue-character}
If a unit $a \in (\ZZ/q\ZZ)^\times$ has $V_N(a) = 0$ (zero visits), then
there exists a nontrivial character~$\chi$ with
$\lVert S_\chi(N) \rVert^2 \geq N^2 / (q{-}2)^2$.
In particular, $\lVert S_\chi(N) \rVert \geq N/(q{-}2)$: the
character sum grows \emph{linearly}, contradicting SVE\@.
\end{theorem}

\begin{theorem}[SVE contradicts avoidance ---
\lean{EM/Population/AvoidanceTube.lean}{sve\_contradicts\_avoidance}
{\color{leanink}$\checkmark$}]
\label{thm:sve-contradicts-avoidance}
SVE $\;\Longrightarrow\;$ for all large~$N$, no unit $a$ can have
$V_N(a) = 0$.
\end{theorem}

\paragraph{Shielding lemma and spectral conspiracy.}
The \defname{shielding lemma} observes that $\seq{2}(n{+}1)$ is always
a term of the sequence, hence never a missing prime:

\begin{theorem}[Shielding lemma ---
\lean{EM/Population/InfiniteM.lean}{shielding\_lemma}
{\color{leanink}$\checkmark$}]
\label{thm:shielding}
For every~$n$, $\seq{2}(n{+}1) \notin M$.
\end{theorem}

\noindent
When a missing prime~$q$ divides $\Prod{2}(n)+1$, some strictly smaller
appeared prime shields~$q$ from capture
(\lean{EM/Population/InfiniteM.lean}{factor\_dichotomy\_strong}
{\color{leanink}$\checkmark$}).
Under~$\neg$WM, this shielding mechanism creates a \emph{spectral conspiracy}:

\begin{theorem}[Spectral conspiracy ---
\lean{EM/Population/SpectralConspiracy.lean}{spectral\_conspiracy\_landscape}
{\color{leanink}$\checkmark$}]
\label{thm:spectral-conspiracy}
$\neg\mathrm{WM}$ implies:
\begin{enumerate}[nosep]
\item $M$ is infinite
  (\lean{EM/Population/SpectralConspiracy.lean}{not\_wm\_implies\_missing\_infinite}
  {\color{leanink}$\checkmark$}).
\item The joint excess energy $\sum_{q \in M} \log\frac{q-1}{q-2}$ diverges.
\item For any finite ``guardian pool'' $G$ of appeared primes, the
  demand $\sum_{q \in F} 1/(q{-}1)$ from a large enough subset
  $F \subseteq M$ exceeds~$|G|$.
\item Arbitrarily many missing primes exist, each witnessing a nontrivial
  character with linear growth (Theorem~\textup{\ref{thm:rogue-character}}).
\end{enumerate}
\end{theorem}

\noindent
The guardian exhaustion clause~(3) formalizes the intuition that finitely
many appeared primes cannot ``explain away'' the avoidance at infinitely
many missing primes.  The three files
\code{EM/Population/AvoidanceTube.lean} (516~lines),
\code{EM/Population/InfiniteM.lean} (303~lines),
\code{EM/Population/SpectralConspiracy.lean} (306~lines)
total 1{,}125~lines, all zero~\code{sorry}.

\paragraph{The shifted squarefree population.}
The EM accumulator $\Prod{2}(n)$ is squarefree
(\lean{EM/Population/WeakErgodicity.lean}{prod\_squarefree}),
so every Euclid number $\Prod{2}(n)+1$ belongs to the shifted
squarefree population
$\mathcal{S} = \{m \geq 2 : m{-}1 \text{ squarefree}\}$
(density $6/\pi^2$).  An earlier draft decomposed the orbit
equidistribution problem accordingly into a \emph{population
equidistribution} (PE: among $q$-rough elements of $\mathcal{S}$, the
smallest prime factor is equidistributed modulo~$q$) and a
\emph{population transfer} (PT: the EM trajectory samples
$\mathcal{S}$ without mod-$q$ bias).  PE is false---the class
densities are head-dominated convergent series (Dead End~\#160,
\S\ref{sec:ant-def})---so the decomposition is retired; the
population $\mathcal{S}$ itself remains the right ambient set for the
ensemble arguments below.

